\chapter{Технологическая часть}

\section{Выбор средств программной реализации}

\subsubsection{Выбор языка программирования}

В качестве языка программирования был выбран Python так как:
\begin{itemize}[label*=---]
	\item  поддерживает библиотеки для анализа данных (pandas);~\cite{pandas}
	\item  обладает инструментами для реализации машинного обучения (pytorch,
	catboost, scikit-learn);~\cite{pytorch, catboost, scl}
	\item обладает большим сообществом разработчиков и является одним из
	наиболее часто используемых языков в сфере машинного обучения.
\end{itemize}

\subsubsection{Среда разработки}

Для разработки программного обеспечения была выбрана интегрированная среда разработки (IDE) Visual Studio Code (VS Code), который обладает следующими преимуществами:
\begin{enumerate}
	\item Производительность и низкое потребление ресурсов.
	VS Code отличается высокой скоростью работы и умеренным использованием системных ресурсов;~\cite{vsc}
	\item Экосистема расширений.
	Магазин расширений (Extensions Marketplace) предоставляет доступ к множеству плагинов для кастомизации среды под конкретные нужды;~\cite{vsc}
	\item Активное сообщество и регулярные обновления.
	VS Code разрабатывается Microsoft при поддержке open-source-сообщества, благодаря чему выходят частые обновления с новыми функциями и исправлениями багов.~\cite{vsc}
\end{enumerate}


\section{Реализация программного обеспечения}

Реализованное программное обеспечение реализует метод дифференциальной приватности при генерации синтетических данных.
Написанная программа включает в себя компоненты, которые организованы в модули с разделением ответственности.

На рисунке~\ref{img:app_struct} представлена диаграмма модулей.
\includeimage
{app_struct} % Имя файла без расширения (файл должен быть расположен в директории inc/img/)
{f} % Обтекание (с обтеканием)
{h!} % Положение рисунка (см. wrapfigure из пакета wrapfig)
{0.8\textwidth} % Ширина рисунка
{Диаграмма модулей} % Подпись рисунка
% \imgs{index-struct}{h!}{1}{Индекс как структура данных}


\section{Основные модули}

\subsubsection{Модуль src}

Модуль src содержит реализацию алгоритма построения Байесовской сети с дифференциальной приватностью и отвечает за построение синтетических данных на основе ее структуры.


Класс Describer реализует методы:
\begin{enumerate}[label*=---]
	\item init --- конструктор класс, к котором инициализируются атрибуты класса;
	\item describe\_dataset --- описывает датасет с использованием байесовской сети;
	\item read\_dataset\_from\_csv --- загружает датасет из CSV-файла;
	\item infer\_attribute\_data\_types --- определяет типы данных атрибутов;
	\item analyze\_dataset\_meta --- анализирует метаинформацию датасета (ключи, категориальность и т.д.);
	\item is\_categorical --- определяет, является ли атрибут категориальным;
	\item represent\_input\_dataset\_by\_columns --- представляет входной датасет в виде объектов-атрибутов;
	\item inject\_laplace\_noise\_into\_distribution\_per\_attribute --- добавляет лапласовский шум в распределения атрибутов;
	\item encode\_dataset\_into\_binning\_indices --- кодирует датасет в индексы бинов для построения байесовской сети.
	\item save\_dataset\_description\_to\_file --- cохраняет описание датасета в файл.
\end{enumerate}


Класс Generator реализует методы:

\begin{enumerate}[label*=---]
	\item init --- онструктор класса. Инициализирует атрибуты экземпляра;
	\item generate\_dataset --- Генерирует синтетический датасет с учетом зависимостей между атрибутами;
	\item save\_synthetic\_data --- Сохраняет сгенерированный синтетический датасет в файл;
	\item get\_sampling\_order --- Возвращает порядок выборки атрибутов на основе структуры байесовской сети;
\end{enumerate}

\subsubsection{Модуль lib}

Модуль lib содержит вспомогательные функции для работы модуля src.
В данном модуле реализованы функции:
\begin{enumerate}[label*=---]
	
	\item calculate\_sensitivity --- вычисляет чувствительность для построения байесовской сети;
	\item calculate\_delta --- вычисляет параметр delta для применения дифференциальной приватности;
	\item calculate\_k --- вычисляет максимальную степень при построении байесовской сети;
	\item greedy\_bayes --- строит байесовскую сеть с помощью жадного алгоритма;
	\item exponential\_mechanism --- реализует экспоненциальный механизм для выбора структуры сети с учетом параметра приватности;

	\item set\_random\_seed --- станавливает фиксированное значение seed для генераторов случайных чисел;
	\item mutual\_information --- вычисляет взаимную информацию между двумя распределениями (Series/DataFrame), используя mutual\_info\_score из sklearn;
	\item pairwise\_attributes\_mutual\_information --- вычисляет нормализованную взаимную информацию для всех пар признаков;
	\item normalize\_given\_distribution --- нормализует переданный массив частот (преобразует в вероятностное распределение) и обрабатывает отрицательные и бесконечные значения;
	\item read\_json\_file --- читает и возвращает содержимое JSON-файла.
\end{enumerate}
\newpage
\section{Реализация алгоритма построения БС с ДП}


Реализация алгоритма для построения структуры Байесовской сети с дифференциальной приватностью представлена на листинге~\ref{lst:gb1.py}.

\includelisting
{gb1.py} % Имя файла с расширением (файл должен быть расположен в директории inc/lst/)
{Реализация алгоритма для построения структуры Байесовской сети с дифференциальной приватностью --- часть 1} % Язык программирования (необязательный аргумент)% Подпись листинга

\includelisting
{gb2.py} % Имя файла с расширением (файл должен быть расположен в директории inc/lst/)
{Реализация алгоритма для построения структуры Байесовской сети с дифференциальной приватностью --- часть 2} % Язык программирования (необязательный аргумент)% Подпись листинга

Реализация алгоритма генерации синтетических данных на основе Байесовской сети представлена на листинге~\ref{lst:generator.py}.

\includelisting
{generator.py} % Имя файла с расширением (файл должен быть расположен в директории inc/lst/)
{Реализация алгоритма генерации синтетических данных на основе Байесовской сети} % Язык программирования (необязательный аргумент)% Подпись листинга

\section{Формат входных и выходных данных}
\subsubsection{Входные данные}
Входными данными для разработанного методы являются:
\begin{enumerate}
	\item датасет с исходными данными --- CSV-файл с названиями атрибутов в формате на листинге~\ref{lst:csv.csv}:
	\item настройки приватности --- числовой параметр $\epsilon$, который учитывется при добавлении шума;
	\item количество генерируемых записей --- числовой параметр, который учитывается непосредственно при генерации синтетических данных;
	\item категориальные атрибуты --- строковый параметр, в котором через запятую перечисляются названия атрибутов, являющихся категориальными;
	\item количество родительских узлов --- параметр, который учитывается при построении Байесовской сети.
\end{enumerate}

\includelisting
{csv.csv} % Имя файла с расширением (файл должен быть расположен в директории inc/lst/)
{CSV-файл с названиями атрибутов} % Подпись листинга

\subsubsection{Выходные данные}

Результатом работы методы являются синтетические данные с обеспечением дифференциальной приватности в формате CSV-файла с заданным количеством записей.

\section{Тестирование программного обеспечения}

\subsubsection{Методика тестирования}

Тестирование реализованного программного обеспечения проводилось в комплексном формате и включало в себя типы тестов: модульное тестирование --- метод тестирования, при котором проверяется корректности работы отдельных компонентов программы в изоляции от остальных частей системы; функциональное тестирование --- метод тестирования, при котором проверяется корректноая работа системы в целом.

\subsubsection{Результаты тестирования}

Результаты тестирования показали корректность работы разработанного метода для различных входных данных. Были протестированы следующие аспекты:
\begin{enumerate}[label*=---]
	\item Загрузка датасета из файла;
	\item Анализизирование датасета;
	\item Создание структуры Байесовской сети;
	\item Добавление шума;
	\item Генерация синтетических данных;
	\item Сохранение выходных данных.
\end{enumerate}

Все тесты были пройдены успешно.

\newpage
\section{Пример пользовательского интерфейса}

Интерфейс пользователя приложения представлен на рисунке~\ref{img:primer}.
В интерфейсе присутствуют поля для ввода данных, также многие поля изначальная заполнены данными по умлчанию.
\includeimage
{primer} % Имя файла без расширения (файл должен быть расположен в директории inc/img/)
{f} % Обтекание (с обтеканием)
{h!} % Положение рисунка (см. wrapfigure из пакета wrapfig)
{0.8\textwidth} % Ширина рисунка
{Реализация пользовательского интерфейса} % Подпись рисунка
% \imgs{index-struct}{h!}{1}{Индекс как структура данных}


% \section{Реализация алгоритма построяения Байесовской сети}




